%%=====================================================================
%% preamble.excerpt.mw.tex
%% Merriam-Webster 版词汇讲义专用导言区。
%%
%% 与 preamble.excerpt.tex（UApiPro 版）并列存在，互不影响。
%% 差别只在数据来源带来的栏目变化：
%%   * 没有「搭配」栏 —— MW 不提供搭配数据，故不再需要 \wcol；
%%   * 新增 \wtag{标签}{内容}，用于词形变化 / 近义词 / 反义词 / 词源
%%     这些「一条标签 + 一段内容」的附注行；
%%   * 义项与例句的关系是 MW 数据自带的，所以 \wsense / \wex 的用法不变。
%%
%% 本文件刻意写成**自包含**的：不 \input 另一份导言区。
%% 因为 TeX 的 \input 是相对编译目录解析的，而 main.tex 与导言区往往不在
%% 同一层，静态文件里写死相对路径很容易在换目录后失效。
%%
%% 版式约定（每个词条）：
%%   1. \whead{词}              词条标题
%%   2. \wsent{…\wmark{词}…}    原文原句（加粗），标记词加下划线
%%   3. \wsense{词性}{英文释义}  一个义项一行；词性标注保留
%%        \wex{例句}            该义项自带的例证（不加粗）
%%      \wtag{词形变化}{…}       附注行
%%   4. \wblank                 非单词标记：释义处留空
%%=====================================================================

\usepackage[]{geometry}
\geometry{a4paper,margin=0.7cm}
\geometry{twoside, bindingoffset=1cm}
\linespread{1.1}
\usepackage[UTF8]{ctex}
\usepackage{fontspec}
\usepackage{xcolor}
% ulem 的 \uline 可以在行尾断开；normalem 保留 \emph 原本的斜体语义
\usepackage[normalem]{ulem}
\usepackage{hyperref}

%%-------------------------------------------------- 字体
\setmainfont{Times New Roman}
\setsansfont{Arial}
\setmonofont{Courier New}
\setCJKmainfont{SimSun}[AutoFakeBold=2.0]
\setCJKsansfont{SimHei}
\pagestyle{plain}

%%-------------------------------------------------- 配色
\definecolor{wInk}{RGB}{40, 40, 40}
\definecolor{wHead}{RGB}{24, 74, 122}
\definecolor{wPos}{RGB}{133, 79, 20}
\definecolor{wDef}{RGB}{53, 94, 59}
\definecolor{wRule}{RGB}{100, 100, 100}

%%-------------------------------------------------- 排版参数
\setlength{\parindent}{0pt}
\setlength{\parskip}{0pt}
% 例句 / 附注行的统一缩进宽度
\newlength{\wIndent}
\setlength{\wIndent}{1.2em}
% 相邻词条之间的间距（拉开相邻的标记词）
\newlength{\wEntryGap}
\setlength{\wEntryGap}{1.6em}

%%-------------------------------------------------- 词条之间的间距
\newcommand{\wgap}{\par\vspace{\wEntryGap}}

%%-------------------------------------------------- 词条标题
% \whead{abandon}
\newcommand{\whead}[1]{%
  \wgap
  {\color{wHead}\bfseries\large #1}%
  \par\nopagebreak\smallskip
}

%%-------------------------------------------------- 原文原句（置顶，加粗）
% \wsent{The committee decided to \wmark{abandon} the proposal.}
\newcommand{\wsent}[1]{%
  \par\nopagebreak
  {\bfseries #1}%
  \par\nopagebreak
  %\vspace{1.2pt}%
  %{\color{wRule}\hrule height 0.6pt}%
  %\par\nopagebreak\vspace{4pt}%
}

%%-------------------------------------------------- 标记词：下划线
% 用 \uline 而非 \underline：整句被标记时也能正常断行
\newcommand{\wmark}[1]{\uline{#1}}

%%-------------------------------------------------- 义项：词性 + 英文释义
% 一个义项写一条。MW 的义项自带例证，紧随其后的 \wex 即属于该义项。
% \wsense{verb}{to give up completely}
\newcommand{\wsense}[2]{%
  \par\nopagebreak\smallskip
  {\color{wPos}\bfseries [#1]}\ {\color{wDef}#2}%
  \par\nopagebreak
}

%%-------------------------------------------------- 例句（不加粗）
% \wex{They abandoned the ship.}
\newcommand{\wex}[1]{%
  \par\nopagebreak
  \hspace*{\wIndent}#1%
  \par\nopagebreak
}

%%-------------------------------------------------- 附注行：标签 + 内容
% 用于词形变化 / 近义词 / 反义词 / 词源 / 首次使用。
% \wtag{词形变化}{past tense: abandoned}
\newcommand{\wtag}[2]{%
  \par\nopagebreak\smallskip
  \hspace*{\wIndent}{\color{wPos}\bfseries #1:\,}\ {\color{wInk}#2}%
  \par\nopagebreak
}

%%-------------------------------------------------- 非单词标记：留空
% 加粗的若不是单个单词（短语、整句等），其释义处留白，供自行补注。
\newcommand{\wblank}{%
  \par\nopagebreak
  \hspace*{\wIndent}\rule{0pt}{2.6em}%
  \par\nopagebreak
}
